\subsection*{Writing the fact requested in a human Chinese question}
Two guests asked about the same Beijing restaurant. One wanted its opening hours. The other asked where it was without using the label for address. Given identical catalogue rows, the visual writer printed the hours in the first reply and the street address in the second (Fig.~\ref{fig:visual-fact-reply}a). On 515 eligible human test turns it returned the requested value without another catalogue value or an extraneous number on 507. Erasing the question reduced this count to 237. Both answers were correct on 111 of 114 pairs of naturally occurring questions about the same card, compared with zero after erasure. When both the question string and venue were absent from fitting, the writer succeeded on 57 of 59 turns. Of 80 new wordings lacking the literal target field name, it answered 78 correctly, compared with 71 for character-trigram retrieval and 28 for matching field names with a domain fallback.

Of the 515 replies, 194 reproduced the full human answer and 476 explicitly named the correct relation and included the clean value. Across outputs there were 136 distinct wordings after masking copied spans. One hours answer opened with a malformed clause, and one price answer added an unsupported judgement. On the first field-selection error the guest asked for hotel type and the writer gave an address (Fig.~\ref{fig:visual-fact-reply}b). Factual slot recovery alone overstates the quality of the complete sentence. The matched byte reader recovered 508 clean values in 104.6 seconds of fitting, compared with 109.8 seconds for the optical reader on the same RTX 4090 D.

The original reader often lost track of which property a guest wanted when two properties occurred in one question. On 54 paired tests of three controlled constructions it got both directions right nine times. A 40,657-parameter selector trained on the visual context of two field mentions raised this count to 54. On twelve further constructions, complete pairs rose from 58 to 180 of 216. Successes occupied ten entire constructions. Every pair failed on two other constructions, one dismissing a field after an initial request and one affirming a request after an initial negation. The natural test score remained 507. These probes measure transfer across sentence constructions under exact printed field matching.

In the separate English writer, learned copy actions carried a pictured restaurant and landmark name into new sentences (Fig.~\ref{fig:visual-fact-reply}c). On Chinese cards complete normalized factual descriptions rose from 1,606 to 1,668 of 1,847. Both directions of a name exchange preserved the original ink and its role in all 128 pairs tested with new Latin names or historical marks. Requiring the entire description to convey its other recognized facts reduced the two counts to 59 and 56 of 64. Historical marks here function as names supplied on the card.
